\section{Implementación del Subsistema de Gestión}

El subsistema de gestión constituye el punto de acceso único al sistema: recibe
las imágenes capturadas en los accesos, coordina su procesamiento y expone la
información resultante. Este apartado documenta su implementación efectiva y
cierra con una declaración explícita del estado de cada componente junto con la
evidencia que lo respalda, con el propósito de distinguir con precisión lo
construido de lo pendiente.

\subsection{Modelo físico de datos y migraciones}

El esquema físico comprende las tres entidades especificadas en el diseño
---vehículo, sesión y detección---, cuyas columnas se detallan en la Tabla
\ref{tab:diccionario_datos}. La implementación añade las siguientes decisiones
que no se derivan directamente del modelo conceptual:

\begin{itemize}
    \item \textbf{Identificadores generados en la aplicación.} Las claves
        primarias son identificadores únicos universales generados por el
        servicio antes de la inserción, y no valores secuenciales asignados por
        el motor de base de datos. Esto permite que el servicio conozca el
        identificador de la sesión en el mismo instante en que la crea, sin
        requerir una lectura posterior, y facilita la correlación de registros
        entre componentes.

    \item \textbf{Procedencia del vehículo preservada.} Cada sesión referencia
        obligatoriamente a un vehículo, y el borrado del vehículo arrastra sus
        sesiones asociadas. La clave foránea es la que garantiza que no exista
        una sesión sin vehículo identificado.

    \item \textbf{Caja delimitadora como documento estructurado.} Las
        coordenadas de las detecciones se almacenan en un campo de tipo documento
        estructurado y no en cuatro columnas escalares. La decisión responde a
        que la caja es un dato de auditoría y no un atributo de consulta: no se
        filtra por coordenadas, pero interesa conservar su estructura para
        reconstruir una detección.

    \item \textbf{Estado de la sesión como enumeración.} El ciclo de vida de una
        sesión se restringe a los valores activa, completada o cancelada, de modo
        que el motor de base de datos rechace cualquier estado no previsto.

    \item \textbf{Marcas temporales con zona horaria.} Todos los instantes se
        almacenan con zona horaria y valor por omisión asignado por el servidor,
        lo que evita ambigüedades cuando el servicio y la base de datos residen en
        contenedores con configuración regional distinta.
\end{itemize}

El modelo declara índices sobre la matrícula, sobre el identificador de vehículo
y el estado de la sesión, y sobre el identificador de sesión y la cámara en la
bitácora de detecciones. Estos índices corresponden a los patrones de consulta
efectivos del sistema: recuperar el historial de un vehículo, listar las sesiones
activas para calcular la ocupación y reconstruir las observaciones de una visita.

El acceso a los datos se resuelve mediante una dependencia que entrega una sesión
asíncrona por solicitud y delimita su transacción: los cambios se confirman al
concluir la atención de la solicitud y se revierten si se produce una excepción.
Esta política garantiza que una operación que involucra varias escrituras
---crear el vehículo y crear su sesión--- se aplique de forma atómica. El motor
mantiene un grupo de conexiones reutilizables, dimensionado para sostener la
concurrencia esperada del servicio sin saturar la base de datos.

El esquema no se crea de forma directa: se materializa mediante tres revisiones
versionadas que se aplican en orden y dejan registro del cambio estructural. La
primera establece las tres entidades con sus restricciones e índices; la segunda
añade el motivo y el operador del cierre manual, que se registran junto al
instante de salida previsto en el esquema inicial; la tercera retira columnas que
ningún componente produce: el vector de características, que se midió y se
descartó como señal de identidad, y los campos de marca, que ningún clasificador
entrega. Esa tercera revisión merece señalarse porque documenta una decisión:
conservar campos que nadie escribe obliga a cada lector del esquema a preguntarse
cuáles son reales.

La aplicación de las tres revisiones sobre una base vacía es reproducible y forma
parte de la verificación automática: el banco de pruebas de integración levanta
una base temporal, aplica las migraciones y ejercita sobre ella el ciclo completo
de registro y cierre.

\subsection{Interfaz de programación}

La interfaz se expone bajo un prefijo versionado y agrupa sus recursos por área
funcional: operación, consulta, autenticación y diagnóstico. El ingreso de
fotogramas no utiliza HTTP: el proceso de captura publica la tarea de inferencia
en la cola y el trabajador escribe el resultado directamente en la base de datos,
de modo que la interfaz no se interpone en el camino crítico. La Tabla
\ref{tab:api_implementada} enumera los recursos efectivamente implementados y sus
parámetros.

\newcolumntype{P}[1]{>{\raggedright\arraybackslash}m{#1\dimexpr\textwidth\relax}}
\begin{longtable}{
        P{0.09}
        P{0.27}
        P{0.34}
        P{0.20}
    }
    \caption{Recursos implementados de la interfaz de programación}
    \label{tab:api_implementada}\\

    \toprule
    \rowcolor{gray!20}
    \textbf{Método} &
    \textbf{Ruta} &
    \textbf{Función} &
    \textbf{Parámetros} \\
    \midrule
    \endfirsthead

    \multicolumn{4}{c}{\textit{Continuación de la Tabla
    \ref{tab:api_implementada}}} \\
    \toprule
    \rowcolor{gray!20}
    \textbf{Método} &
    \textbf{Ruta} &
    \textbf{Función} &
    \textbf{Parámetros} \\
    \midrule
    \endhead

    GET & /api/v1/stats/active & Lista los vehículos actualmente dentro del
    recinto, con sus atributos. & — \\
    \midrule

    GET & /api/v1/stats/history & Lista las sesiones completadas, ordenadas de la
    más reciente a la más antigua. & Límite ($1$--$500$) y desplazamiento. \\
    \midrule

    GET & /api/v1/stats/occupancy & Devuelve el número de sesiones activas y el
    instante de la consulta. & — \\
    \midrule

    GET & /api/v1/stats/daily & Devuelve, por día, el total de sesiones y la
    duración media. & Número de días ($1$--$90$). \\
    \midrule

    WS & /api/v1/ws/occupancy & Canal persistente que propaga las variaciones de
    ocupación. & — \\
    \midrule

    POST & /api/v1/auth/login & Autentica al operador y entrega el token de
    sesión. & Usuario y contraseña. \\
    \midrule

    POST & /api/v1/auth/logout & Revoca la sesión vigente. & — \\
    \midrule

    GET & /api/v1/auth/me & Devuelve el operador autenticado. & — \\
    \midrule

    GET & /api/v1/cameras & Estado y último latido de cada cámara. & — \\
    \midrule

    POST & /api/v1/sessions/\{id\}/close & Cierra una sesión activa con operador
    y motivo, y deja registro de auditoría. & Identificador y motivo. \\
    \midrule

    GET & /api/v1/vehicles/search & Busca vehículos por matrícula o atributos.
    & Texto de búsqueda. \\
    \midrule

    GET & /api/v1/vehicles/\{id\} & Detalle del vehículo con su historial de
    sesiones. & Identificador. \\
    \midrule

    GET & /api/v1/vehicles/\{id\}/image & Entrega el recorte asociado al
    vehículo. & Identificador. \\
    \midrule

    GET & /api/v1/stats/model & Distribución de tipos y colores observados.
    & — \\
    \midrule

    GET & /api/v1/stats/colors & Agregado de sesiones por color predominante.
    & — \\
    \midrule

    GET & /health, /health/live y /health/ready & Disponibilidad y preparación
    del servicio para recibir tráfico. & — \\*
    \multicolumn{4}{r}{\footnotesize\textit{Fuente: Elaboración propia}} \\*
    \bottomrule
\end{longtable}

El recurso de historial implementa paginación por límite y desplazamiento ---con
un máximo de quinientos registros por solicitud--- y el de agregados diarios
acota el intervalo consultable, evitando que una petición sin restricciones
recorra la totalidad de la bitácora. Los contratos de entrada y salida se
declaran como tipos validados, de modo que una solicitud malformada se rechace
antes de alcanzar la lógica de negocio. La especificación formal de la interfaz
se publica de forma automática a partir de esas declaraciones, lo que garantiza
que la documentación no se desincronice de la implementación.

El acceso está protegido: los recursos de lectura no responden sin un token de
sesión válido, que el cliente conserva en una cookie inaccesible al código del
navegador. La distinción entre recursos de escritura ---destinados a los
dispositivos de captura--- y de consulta ---destinados al personal del
parqueadero--- se mantiene en la autorización y no solo en la nomenclatura de las
rutas.

El cierre de una sesión admite dos caminos. El primero es automático: cuando el
reconocedor entrega una matrícula, la política de emparejamiento busca entre las
sesiones activas una coincidencia exacta de matrícula, exige que ambas lecturas
---la del ingreso y la del egreso--- superen un umbral de confianza de $0{,}6$ y,
si varias sesiones comparten la matrícula, cierra la más reciente. Si no hay
coincidencia o la confianza es insuficiente, \textbf{no devuelve la mejor
aproximación}: deja la sesión abierta. La decisión es deliberada y se midió antes
de adoptarla. Se evaluó primero la alternativa barata, usar el extractor del
clasificador de tipo como vector de identidad, sobre $400$ recortes del conjunto
congelado: la similitud promedio entre el mismo vehículo fue $0{,}860$, mientras
que entre vehículos distintos el percentil $90$ alcanzó $0{,}973$ y el $99$
llegó a $0{,}991$. Es decir, el $19{,}9\%$ de los pares de vehículos diferentes
supera el umbral de $0{,}95$, de modo que ningún punto de corte los separa: el
extractor está entrenado para la silueta, no para la identidad. El segundo camino
es el cierre manual por parte del operador, con motivo y responsable registrados
en los campos de auditoría. La regla que gobierna ambos caminos es la misma: un
egreso sin emparejar es visible en el nivel de ocupación, mientras que un
emparejamiento incorrecto vacía el recinto de forma silenciosa.

\subsection{Procesamiento asíncrono}

La aplicación de tareas se configura con la cola como intermediario de mensajes y
como almacén de resultados, serialización en formato de texto estructurado y
tiempo universal coordinado como zona horaria de referencia. Se declaran dos
tareas, correspondientes al procesamiento de una imagen de ingreso y de una
imagen de salida, y se impone un límite temporal a su ejecución: cinco minutos
como máximo absoluto y cuatro como límite flexible, de modo que una tarea
bloqueada no ocupe indefinidamente al trabajador. El registro del inicio de cada
tarea se encuentra habilitado, lo que permite distinguir una tarea encolada de
una en ejecución.

Las dos tareas se encuentran implementadas y su cuerpo está separado de la cola
para poder ejercitarlo sin intermediario. La tarea decodifica el fotograma,
interpreta el instante de captura y delega en un manejador que concentra la
lógica: invoca al motor de visión y mide la latencia de esa inferencia, agrega las
detecciones por trayectoria hasta confirmar un evento, persiste el recorte,
invoca al lector de placas sobre la caja del vehículo y registra la observación de
forma atómica junto con la publicación del evento. El instante que se persiste es
el de muestreo del fotograma y no el de ejecución de la tarea: un atraso en la
cola no debe desplazar un vehículo a una hora posterior. Un fotograma que no
puede decodificarse se informa y no produce efecto alguno. La agregación por
trayectoria garantiza la idempotencia del camino de escritura: tres fotogramas de
un mismo vehículo producen un vehículo y una sesión, no tres.

\subsection{Notificación en tiempo real}

El canal de notificación se implementa mediante un gestor de conexiones que
mantiene la lista de clientes suscritos, incorpora una conexión al aceptarla,
la retira al cerrarse y difunde un mensaje a todos los suscriptores. La difusión
es tolerante a fallos: si el envío a un cliente produce una excepción, esa
conexión se marca como inactiva y se retira del conjunto, de manera que un
cliente desconectado de forma abrupta no interrumpa la notificación al resto.

Los mensajes se estructuran por tipo, distinguiendo el registro de un ingreso, el
de una salida, la actualización del nivel de ocupación y un mensaje de
mantenimiento del canal. Esta clasificación permite que el cliente actualice
únicamente la porción de la interfaz afectada en lugar de recargar el estado
completo.

\subsection{Aplicación web}

El cliente se construye con React y TypeScript sobre Vite, y organiza su
navegación en un componente de disposición común que envuelve las vistas y
redirige al panel cualquier ruta no reconocida. Se implementan seis vistas: el
acceso con autenticación, el panel de monitoreo con el estado del recinto y el
flujo de eventos, el historial de sesiones, la consulta de incidentes, el padrón
de vehículos con su detalle e historial, y la vista de desempeño del modelo con la
distribución de tipos y colores observados. El acceso a datos se centraliza en un
servicio que declara tipos para cada respuesta de la interfaz de programación, de
modo que el compilador detecte una discrepancia entre el contrato esperado y el
recibido.

La actualización en tiempo real se encapsula en un enlace que establece la
conexión con el canal persistente al montar el componente y la libera al
desmontarlo, exponiendo el último mensaje recibido. Se implementa además un
mecanismo de consulta periódica como alternativa para los datos que no dependen
de la notificación, lo que permite que la interfaz permanezca operativa aun si el
canal persistente no está disponible.

\subsection{Estado de implementación del subsistema}
\label{sec:estado_gestion}

La Tabla \ref{tab:estado_gestion} declara el estado de cada componente del
subsistema de gestión y la evidencia que lo respalda. Se explicita con el
propósito de que la evaluación posterior permita atribuir cada resultado al
componente que efectivamente lo produjo.

\begin{longtable}{
        P{0.30}
        P{0.15}
        P{0.47}
    }
    \caption{Estado de implementación del subsistema de gestión}
    \label{tab:estado_gestion}\\

    \toprule
    \rowcolor{gray!20}
    \textbf{Componente} &
    \textbf{Estado} &
    \textbf{Observación} \\
    \midrule
    \endfirsthead

    \multicolumn{3}{c}{\textit{Continuación de la Tabla
    \ref{tab:estado_gestion}}} \\
    \toprule
    \rowcolor{gray!20}
    \textbf{Componente} &
    \textbf{Estado} &
    \textbf{Observación} \\
    \midrule
    \endhead

    Modelo de datos y esquema & Implementado & Las tres entidades se materializan
    con sus restricciones e índices. \\
    \midrule
    Versionado de migraciones & Implementado & Tres revisiones aplicadas en orden;
    la aplicación sobre una base vacía es reproducible y se ejercita en la
    verificación automática. \\
    \midrule
    Registro de ingreso & Implementado & Crea el vehículo y la sesión, y responde
    de inmediato. Verificado en pruebas de caso de uso y de integración. \\
    \midrule
    Registro de salida & Implementado & Emparejamiento por matrícula exacta con
    umbral de confianza, y cierre manual auditado como alternativa. \\
    \midrule
    Control de autorización & Pendiente & El contraste contra un padrón
    (RF-10) requiere una entidad de vehículos habilitados; la entidad actual
    registra vehículos observados en los accesos. \\
    \midrule
    Recursos de consulta & Implementado & Ocupación, sesiones activas, historial
    paginado, agregados diarios, padrón, detalle de vehículo y distribución de
    atributos. \\
    \midrule
    Tareas de inferencia & Implementado & La tarea decodifica el fotograma, invoca
    al motor de visión, agrega por trayectoria y persiste los eventos
    confirmados con su latencia. \\
    \midrule
    Notificación en tiempo real & Implementado & Canal persistente con difusión
    tolerante a desconexiones. \\
    \midrule
    Aplicación web & Implementado & Seis vistas: acceso, panel, historial,
    incidentes, padrón con detalle y desempeño del modelo. \\
    \midrule
    Pruebas automatizadas & Implementado & $140$ pruebas unitarias y $4$ de
    integración contra una base real con las migraciones aplicadas; cobertura
    del $88\%$ sobre el total del código de la aplicación. \\*
    \multicolumn{3}{r}{\footnotesize\textit{Fuente: Elaboración propia}} \\*
    \bottomrule
\end{longtable}
